\documentclass{article}
\usepackage{amsmath,amssymb,amsthm}
\usepackage{algorithm}
\usepackage{algorithmic}
\usepackage{hyperref}
\usepackage{geometry}
\geometry{margin=1in}

\begin{document}

\title{Comprehensive Exposition on Fixed‐Point Differentiation and Optimization}
\author{}
\date{}
\maketitle

\section*{1. Problem Statement}
We consider a map
\[
  F\colon \mathbb{R}^n \times \mathbb{R}^d \;\longrightarrow\; \mathbb{R}^n,
  \quad (x,\theta)\mapsto F(x,\theta),
\]
satisfying:
\begin{enumerate}
  \item[(i)] (\emph{Contraction in \(x\)})\quad
    There exists \(0\le\gamma<1\) such that for all \(x_1,x_2\in\mathbb{R}^n\) and \(\theta\in\mathbb{R}^d\),
    \[
      \|F(x_1,\theta)-F(x_2,\theta)\|\;\le\;\gamma\,\|x_1-x_2\|.
    \]
  \item[(ii)] (\(\mathcal C^1\) smoothness)\quad
    \(F\) is continuously differentiable in both arguments.
\end{enumerate}
By the Banach fixed‐point theorem, for each \(\theta\) there is a unique
\(\Pi(\theta)\in\mathbb{R}^n\) such that
\[
  \Pi(\theta) \;=\; F\bigl(\Pi(\theta),\,\theta\bigr).
\]

\section*{2. Differentiating the Fixed Point}
\subsection*{2.1 Implicit‐Function Formulation}
Define 
\[
  G(x,\theta) := x - F(x,\theta).
\]
Then \(G(\Pi(\theta),\theta)=0\) and
\[
  D_xG = I_n - D_xF(\Pi(\theta),\theta), 
  \quad
  D_\theta G = -\,D_\theta F(\Pi(\theta),\theta).
\]
Since \(\|D_xF\|\le\gamma<1\), the matrix \(I - D_xF\) is invertible.  By the Implicit Function Theorem,
\[
  D_\theta\Pi(\theta)
  \;=\;
  -\bigl(D_xG\bigr)^{-1}D_\theta G
  \;=\;
  \bigl[I_n - D_xF(\Pi(\theta),\theta)\bigr]^{-1}
  \;D_\theta F(\Pi(\theta),\theta).
\]

\subsection*{2.2 Neumann‐Series Representation}
Because \(\|D_xF\|<1\), one may write
\[
  \bigl(I - A\bigr)^{-1}
  = \sum_{k=0}^\infty A^k,
  \quad A := D_xF(\Pi(\theta),\theta),
\]
avoiding explicit inversion in practical algorithms.

\section*{3. Gradient of the Composite Objective}
Let \(S\colon\mathbb{R}^n\to\mathbb{R}\) be \(\mathcal C^1\) and define
\[
  L(\theta) \;=\; S\bigl(\Pi(\theta)\bigr).
\]
Set \(x_\star = \Pi(\theta)\).  Then by the chain rule and the above derivative,
\[
  \nabla_\theta L
  = D_\theta \Pi(\theta)^\top\,\nabla_x S(x_\star)
  = \Bigl[D_\theta F(x_\star,\theta)\Bigr]^\top
    \underbrace{\bigl[I - D_xF(x_\star,\theta)\bigr]^{-\!T}
                   \nabla_x S(x_\star)}_{v},
\]
where \(v\in\mathbb{R}^n\) is the solution of the linear system
\[
  \bigl(I - D_xF(x_\star,\theta)\bigr)^{T}v
  = \nabla_x S(x_\star).
\]

\section*{4. Numerical Algorithms}
\subsection*{4.1 Fixed‐Point Iteration (Forward Pass)}
\begin{algorithm}[H]
\caption{Compute \(x_\star=\Pi(\theta)\) by simple iteration}
\begin{algorithmic}[1]
  \STATE Initialize \(x^{(0)}\), tolerance \(\varepsilon\)
  \FOR{$k=0,1,2,\dots$}
    \STATE \(x^{(k+1)} \leftarrow F\bigl(x^{(k)},\theta\bigr)\)
    \IF{$\|x^{(k+1)}-x^{(k)}\|<\varepsilon$}
      \STATE \textbf{break}
    \ENDIF
  \ENDFOR
  \RETURN \(x_\star := x^{(k+1)}\)
\end{algorithmic}
\end{algorithm}

\subsection*{4.2 Backward Pass: Adjoint Solve}
\begin{itemize}
  \item Compute \(\;g_x = \nabla_x S(x_\star)\).
  \item Solve \((I - J_xF)^{T}v = g_x\) either by:
    \begin{itemize}
      \item \textbf{Conjugate Gradient (CG):} uses only products
        \((I - J_xF)^{T}w\) via vector–Jacobian products.
      \item \textbf{Neumann‐series iteration:}
        \(\;v\leftarrow g_x + D_xF^\top\,v,\) which converges since \(\|D_xF\|<1\).
    \end{itemize}
  \item Finally, \(\;g_\theta = D_\theta F(x_\star,\theta)^{T}\,v\).
\end{itemize}

\subsection*{4.3 Combined Implicit‐Gradient Algorithm}
\begin{algorithm}[H]
\caption{Compute \(\nabla_\theta S(\Pi(\theta))\)}
\begin{algorithmic}[1]
  \STATE \(x_\star \leftarrow\) fixed‐point iteration
  \STATE \(g_x \leftarrow \nabla_x S(x_\star)\)
  \STATE Solve \((I - J_xF)^{T}v = g_x\) by CG or Neumann series
  \STATE \(g_\theta \leftarrow D_\theta F(x_\star,\theta)^{T}\,v\)
  \RETURN \(g_\theta\)
\end{algorithmic}
\end{algorithm}

\section*{5. PyTorch Implementation Sketch}
\begin{verbatim}
import torch
from torch.autograd.functional import vjp

class FixedPointFunction(torch.autograd.Function):
    @staticmethod
    def forward(ctx, theta, F, S, max_iter=50, tol=1e-6):
        # 1) compute fixed point x_star
        x = torch.zeros(F.output_dim)  # or other init
        for _ in range(max_iter):
            x_next = F(x, theta)
            if (x_next - x).norm() < tol:
                break
            x = x_next
        ctx.save_for_backward(x, theta)
        ctx.F, ctx.S = F, S
        return S(x)

    @staticmethod
    def backward(ctx, grad_out):
        x_star, theta = ctx.saved_tensors
        F, S = ctx.F, ctx.S

        # a) nabla_x S(x_star)
        g_x = torch.autograd.grad(S(x_star), x_star, create_graph=True)[0]

        # b) define J_xF^T v via vjp
        def Jt(v):
            _, vjp_fun = vjp(lambda x: F(x, theta), x_star)
            return vjp_fun(v)[0]

        # c) solve (I - J_xF)^T v = g_x via CG
        v = conjugate_gradient(lambda w: w - Jt(w), g_x)

        # d) compute g_theta = D_theta F(x_star,theta)^T v
        _, vjp_theta = vjp(lambda th: F(x_star, th), theta)
        g_theta = vjp_theta(v)[0]

        return grad_out * g_theta, None, None, None, None
\end{verbatim}

\section*{6. Newton and Quasi‐Newton Extensions}
\begin{itemize}
  \item \textbf{Hessian‐vector products:} 
    \[
      H_L(\theta)\,p
      = \frac{d}{d\varepsilon}\Bigl[\nabla L(\theta+\varepsilon p)\Bigr]
      \big|_{\varepsilon=0},
    \]
    via \texttt{torch.autograd.functional.hvp}.
  \item Solve \(H p = -\,\nabla L\) by CG (matrix‐free) or apply L‐BFGS.
  \item Include damping or line‐search strategies as customary.
\end{itemize}

\section*{7. Practical Remarks and Helper Libraries}
\begin{itemize}
  \item Memory complexity remains \(O(d)\), since only \(x_\star\) is stored.
  \item Contraction constant \(\gamma\) governs CG convergence speed; typically
    \(10\text{--}20\) CG iterations suffice if \(\gamma<0.9\).
  \item \textbf{Helper libraries:}
    \begin{itemize}
      \item \texttt{TorchOpt} for \texttt{implicit\_diff} convenience wrappers.
      \item \texttt{DeepEquilibrium.jl}/\texttt{DEQ} layers with Anderson/Broyden solvers.
      \item \texttt{torch.func.fixed\_point} and \texttt{torch.func.vjp} in PyTorch 2.x.
    \end{itemize}
\end{itemize}

\end{document}
